\documentclass[otchet]{../SCWorks}
% Тип обучения (одно из значений):
%    bachelor   - бакалавриат (по умолчанию)
%    spec       - специальность
%    master     - магистратура
% Форма обучения (одно из значений):
%    och        - очное (по умолчанию)
%    zaoch      - заочное
% Тип работы (одно из значений):
%    coursework - курсовая работа (по умолчанию)
%    referat    - реферат
%  * otchet     - универсальный отчет
%  * nirjournal - журнал НИР
%  * digital    - итоговая работа для цифровой кафедры
%    diploma    - дипломная работа
%    pract      - отчет о научно-исследовательской работе
%    autoref    - автореферат выпускной работы
%    assignment - задание на выпускную квалификационную работу
%    review     - отзыв руководителя
%    critique   - рецензия на выпускную работу

% * Добавлены вручную. За вопросами к @mchernigin
\usepackage{../preamble}

\begin{document}

% Кафедра (в родительном падеже)
\chair{математической кибернетики и компьютерных наук}

% Тема работы
\title{Анализ сложности для сортировок, не использующих сравнение элементов}

% Курс
\course{2}

% Группа
\group{251}

% Факультет (в родительном падеже) (по умолчанию "факультета КНиИТ")
\department{факультета компьютерных наук и информационных технологий}

% Специальность/направление код - наименование
% \napravlenie{02.03.02 "--- Фундаментальная информатика и информационные технологии}
% \napravlenie{02.03.01 "--- Математическое обеспечение и администрирование информационных систем}
% \napravlenie{09.03.01 "--- Информатика и вычислительная техника}
\napravlenie{09.03.04 "--- Программная инженерия}
% \napravlenie{10.05.01 "--- Компьютерная безопасность}

% Для студентки. Для работы студента следующая команда не нужна.
\studenttitle{студентки}

% Фамилия, имя, отчество в родительном падеже
\author{Потапкиной Маргариты Андреевны}

% Заведующий кафедрой 
\chtitle{доцент, к.\,ф.-м.\,н.}
\chname{С.\,В.\,Миронов}

% Руководитель ДПП ПП для цифровой кафедры (перекрывает заведующего кафедры)
% \chpretitle{
%     заведующий кафедрой математических основ информатики и олимпиадного\\
%     программирования на базе МАОУ <<Ф"=Т лицей №1>>
% }
% \chtitle{г. Саратов, к.\,ф.-м.\,н., доцент}
% \chname{Кондратова\, Ю.\,Н.}

% Научный руководитель (для реферата преподаватель проверяющий работу)
\satitle{старший преподаватель} %должность, степень, звание
\saname{М.\,И.\,Сафрончик}

% Руководитель практики от организации (руководитель для цифровой кафедры)
\patitle{доцент, к.\,ф.-м.\,н.}
\paname{И.\,А.\,Батраева}

% Руководитель НИР
\nirtitle{доцент, к.\,п.\,н.} % степень, звание
\nirname{И.\,А.\,Батраева}

% Семестр (только для практики, для остальных типов работ не используется)
\term{2}

% Наименование практики (только для практики, для остальных типов работ не
% используется)
\practtype{учебная}

% Продолжительность практики (количество недель) (только для практики, для
% остальных типов работ не используется)
\duration{2}

% Даты начала и окончания практики (только для практики, для остальных типов
% работ не используется)
\practStart{01.07.2024}
\practFinish{13.01.2024}

% Год выполнения отчета
\date{2026}

\maketitle

% Включение нумерации рисунков, формул и таблиц по разделам (по умолчанию -
% нумерация сквозная) (допускается оба вида нумерации)
\secNumbering

\tableofcontents

% Раздел "Обозначения и сокращения". Может отсутствовать в работе
% \abbreviations
% \begin{description}
%     \item ... "--- ...
%     \item ... "--- ...
% \end{description}

% Раздел "Определения". Может отсутствовать в работе
% \definitions

% Раздел "Определения, обозначения и сокращения". Может отсутствовать в работе.
% Если присутствует, то заменяет собой разделы "Обозначения и сокращения" и
% "Определения"
% \defabbr

\intro
В данном отчёте производится анализ сложности и описание особенностей алгоритмов сортировки, не использующих сравнение элементов.

\section{Анализ асимптотики}

\subsection{Сортировка подсчётом}
\small
\inputminted[breaklines]{cpp}{./counting_sort.cpp}
\normalsize

Основные части кода, влияющие на асимптотику алгоритма:

\begin{itemize}
	\item поиск минимума и максимума в массиве для определения диапазона его элементов "--- $O(n)$;
	\item заполнение массива частот "--- $O(n)$;
	\item построение результирующего массива (отсортированного) "--- $O(k + n)$ (внешний цикл работает за $O(k)$, внутренний "--- за $O(n)$ суммарно по всем итерациям внешнего, т.к. сумма чисел в массиве частот равна количеству элементов в изначальном массиве).
\end{itemize}

\textbf{Итого асимптотику алгоритма можно оценить как $O(n + k)$.}

Здесь:
\begin{itemize}
	\item[$\bullet$] $k$ "--- размер диапазона элементов массива ($k = max(a) - min(a)$);
	\item[$\bullet$] $n$ "--- количество элементов в исходном массиве.
\end{itemize}

\subsection{Поразрядная сортировка}
\small
\inputminted[breaklines]{cpp}{./radix_sort.cpp}
\normalsize

Основные части кода, влияющие на асимптотику алгоритма:

\begin{itemize}
	\item внешний цикл по количеству разрядов "--- $O(len)$;
	\item очистка массива частот (для сортировки подсчётом по каждому разряду) "--- $O(b)$;
	\item заполнение массива частот "--- $O(n)$;
	\item построение префиксных сумм по массиву частот "--- $O(b)$;
	\item перезапись массива в отсортированном порядке "--- $O(n)$.
\end{itemize}

\textbf{Итого асимптотику алгоритма можно оценить как $O(len \cdot (n + b))$.}

Здесь:
\begin{itemize}
	\item[$\bullet$] $len$ "--- максимальное количество разрядов в числе;
	\item[$\bullet$] $b$ "--- основание системы счисления;
	\item[$\bullet$] $n$ "--- длина исходного массива.
\end{itemize}

\section{Описание особенностей}

В сортировках, основанных на сравнениях, невозможно достичь асимптотики, лучшей чем $O(n log n)$. Описанные же здесь сортировки не используют сравнения элементов, а потому работают за линейное время.

\subsection{Сортировка подсчётом}

Сортировка подсчётом является стабильной, это проявляется при сортировке сложных структур. Поле, по которому сортируются структуры, должно быть целым числом, при этом оно может быть отрицательным, для этого при обращении к частотному массиву мы выполняем сдвиг на значение минимума, т.е. минимум (который может быть отрицательным) преобразуется в $0$, а максимум "--- в $max(a) - min(a)$.

Асимптотика сортировки подсчётом линейно зависит не только от объёма входных данных, но и от размера диапазона, в котором лежат числа. Соответственно, чтобы, к примеру, отсортировать всего лишь два числа, одно из которых равно $INT\_MAX$, а другое "--- $INT\_MIN$, потребуется затратить порядка $O(INT\_MAX)$ операций, что значительно больше, чем любая, даже самая неэффективная сортировка, основанная на сравнениях. На практике в данном случае сортировка подсчётом попросту неприменима, поскольку мы не можем выделить массив частот такого большого размера. В то же время, сортировка подсчётом применима даже для достаточно больших массивов, если числа в них не слишком велики.

\subsection{Поразрядная сортировка}

Поразрядная сортировка применяется для чисел и строк (здесь в качестве разрядов выступают позиции в строке).

Существует две разновидности поразрядной сортировки: LSD (least significant digit) и MSD (most significant digit). В данном случае рассматривается реализация LSD, которая работает только с натуральными числами. Асимптотика сортировки зависит от количества разрядов в числе, а также от основания системы счисления. Соответственно, максимальное количество разрядов равно логарифму по основанию системы счисления от максимального числа. Соответственно, необходимо грамотно выбрать систему счисления, чтобы минимизировать произведение $log(max(a), base) \cdot base$.

Внутри поразрядной сортировки используется некоторая сортировка, которая должна быть стабильной: это свойство нужно, чтобы при сортировке по старшим разрядам не нарушался относительный порядок по младшим разрядам. В данном случае эффективно использовать сортировку подсчётом, чтобы сохранить линейное время работы.

% После введения — серии \section, \subsection и т.д.

% \conclusion

% Отобразить все источники. Даже те, на которые нет ссылок.
% \nocite{*}

\bibliographystyle{ugost2008}
\bibliography{thesis}

% Окончание основного документа и начало приложений Каждая последующая секция
% документа будет являться приложением
\appendix

\end{document}
